\scanpage{179}
\uline{Proposition} \quad $S$ un préschéma \uline{connexe},

$S'$ un préschéma sur $S$ tel que $S'\longrightarrow S$ soit \uline{submersif} [p. ex. fid$^{t}$ plat et quasi-compact, ou propre surjectif].

$S'',S'''$ les \unsure{carré|carrés}, cube fibré.

On suppose que $S'$ est \uline{somme} de \uline{préschémas} \uline{connexes} (i.e ses composantes connexes sont ouvertes, i.e. $\pi_0(S')$ discret), \deleted{et l'ens des composantes connexes de $S''$ \unsure{compac}}

Soit $\xi$ un pt [géométrique] de $S$. Alors pour un groupe discret $G$ variable, le foncteur
\[
G\leadsto\pi^1(S'/S,\xi\,;G)
\]
\deleted{\unsure{contrav|contravariant}} [où $\pi^1(S'/S,\xi\,;G)$ désigne l'ens des classes, à isom près, de données de descente sur \inserted{le fibré principal trivial} $G_{S'}$ relativement à $S'\to S$, pointées sur $\xi$] est strict.$^{t}$ proreprésentable :
\[
\pi^1(S'/S,\xi\,;G)=\varinjlim\operatorname{Hom}(\pi_i,G)
\]

Si $\pi_0(S'')$ est \inserted{quasi-}\uline{compact}, i.e. \unsure{tte|les} fonction loc constante sur $\pi_0(S'')$ n'a qu'un nb fini de valeurs, -- p. ex. $S''$ quasi-compact], \deleted{\unsure{ou} $S''$ n'a qu'un nb fini de composantes connexes]} \uline{alors} les $\pi_i$ sont de type fini. Si $S''$ a un nb fini \inserted{$n''$} de composantes connexes, alors le foncteur précédent est représentable : \deleted{de \unsure{façon}}
\[
\deleted{\text{\unsure{par}}}\pi^1(S'/S,\xi\,;G)\simeq\operatorname{Hom}(\pi_1(S'/S,\xi),G)
\]
avec $\pi_1(S'/S,\xi)$ \deleted{\unsure{engendré}} \inserted{\unsure{gpe|groupe} discret} de type fini. De façon précise, \deleted{\unsure{si}} le nb $n'$ des comp. conn. de $S'$ est alors fini, et $\deleted{n''-n'}\pi_1(S'/S;\xi)$ est engendré par $n''-n'$ générateurs.

% archive sheet number: 178
